\section{C++20 后端优化总路线：从 reference 到可验收的 CPU/GPU plan}

\subsection{后端优化不能从“写快代码”开始}

很多人第一次写推理后端，会直接跳到 SIMD intrinsic、线程池或者 GPU kernel。这样很容易得到一段看起来很专业、却无法证明正确和无法稳定复用的代码。真正的后端优化要先回答四个问题：

\begin{enumerate}
  \item 正确答案由谁定义。
  \item 当前 baseline 慢在哪里。
  \item 本次优化改变了哪条地址流、生命周期、指令流或同步边界。
  \item oracle 和 profiler 是否同时证明“没算错”和“真的快了”。
\end{enumerate}

这和编译器后端的纪律一样。后端不是随便生成几条目标指令，而是把已经验证的 IR 降到目标机器，并保留语义。AI 推理后端也不是随便调用 kernel，而是把 tensor descriptor、packed weight、memory plan、KV cache 和 runtime segment 落到 CPU 或 GPU 的物理执行。

\begin{keyidea}
后端优化的起点是 reference 和证据，不是 SIMD 或 GPU。没有 reference，就不知道优化后是否正确；没有 profiler，就不知道优化是否击中了真正瓶颈；没有清晰 descriptor，就无法把优化接回编译器和运行时。
\end{keyidea}

\subsection{C++20 后端的基本模块}

一个可维护的 C++20 后端不应该只有 \code{kernel.cpp}。至少要把下面几类职责分开：

\begin{lstlisting}[numbers=none,caption={C++20 后端模块草图}]
backend/cpu/
  cpu_backend.hpp          // backend facade, capability query
  tensor_view.hpp          // non-owning typed/byte views
  aligned_buffer.hpp       // RAII aligned allocation
  packed_weight.hpp        // packed descriptors and cache
  matmul/
    reference_kernel.hpp
    scalar_kernel.hpp
    simd_kernel.hpp
    packing.hpp
  schedule/
    thread_pool.hpp
    work_partition.hpp
  profile/
    cpu_events.hpp

backend/gpu/
  gpu_backend.hpp
  device_buffer.hpp
  stream.hpp
  kernel_adapter.hpp
  workspace.hpp
  staging.hpp
\end{lstlisting}

这些文件名只是示意，真正重要的是边界。C++20 的 RAII 对象负责释放资源，\code{std::span} 或类似轻量视图负责观察连续内存，descriptor 负责描述 shape、stride、dtype、layout 和 quantization，backend facade 负责能力查询和 plan 创建，kernel 只接收已经验证过的输入。

\topic{不要让 kernel 解析模型名。} kernel 应该知道输入指针、shape、stride、scale 和 packed descriptor，不应该知道 \code{layers.12.attn.wq.weight} 这种模型方言。模型名解析属于前端，后端只消费 canonical descriptor。

\topic{不要让 runtime 猜 SIMD 宽度。} runtime 选择 executable plan，但不应该在每个 token 上重新判断 AVX2 或 AVX-512 的细节。CPU backend 初始化时查询能力，编译/prepare 阶段生成候选计划，runtime 只做便宜选择。

\topic{不要让 packed weight 失去来源。} packed buffer 是物理布局，但诊断仍要能回到原始 tensor name、文件 offset、shape 和 checksum。否则一旦输出异常，只能在一坨重排后的字节中调试。

\subsection{CPU 优化阶梯：每一级改变一个机器事实}

CPU 后端可以按阶梯推进，每一级都应该能单独 benchmark 和回滚。

\begin{center}
\begin{tabular}{p{0.17\linewidth}p{0.32\linewidth}p{0.34\linewidth}}
\toprule
阶段 & 改变的机器事实 & 验收证据 \\
\midrule
reference & 定义语义，不追求速度 & oracle 作为比较基准 \\
scalar baseline & 连续循环和清楚地址公式 & 结果正确，profiler 有基线 \\
layout 修正 & stride、对齐、尾块、连续性 & cache miss 和边界路径可解释 \\
packing & 把权重重排成 kernel 友好流 & packing 成本与热路径收益分开统计 \\
SIMD & 向量 load、FMA/dot、累加器展开 & 指令形态和误差边界可验证 \\
线程划分 & work partition、barrier、false sharing & 单线程到多线程 scaling 曲线 \\
NUMA/TLB & 内存放置、页、亲和、预取 & 长 prompt 和多请求下尾延迟稳定 \\
\bottomrule
\end{tabular}
\end{center}

这个阶梯的关键是一次只改变一个主要变量。若同一版同时改 packing、SIMD、线程池和 arena，很难判断速度来自哪里，也很难定位错误。教材后续写具体 kernel 时，每个优化都要说明“输入 descriptor 没变，输出语义没变，只改变了某条物理地址流或执行划分”。

\subsection{SIMD micro-kernel 要写清楚的合同}

SIMD 不是把标量循环换成 intrinsic 就结束。一个合格的 micro-kernel 至少要写清楚：

\begin{itemize}
  \item 支持的 dtype：fp32、fp16/bf16、int8、int4 解包后计算，还是整数 dot。
  \item 输入 layout：行主序、列主序、packed panel、group scale 相对位置。
  \item 向量宽度：AVX2、AVX-512 或运行时选择的抽象宽度。
  \item 累加器数量：同时维护几个输出元素，是否避免 dependency chain。
  \item 尾部路径：\code{K}、\code{N}、\code{head_dim} 不能整除向量宽度时怎样处理。
  \item 数值合同：累加类型、舍入、requantize、允许误差和 reference 比较方式。
  \item 内存合同：输入/输出是否要求对齐，是否允许非连续 stride，是否完整覆盖输出。
\end{itemize}

\begin{lstlisting}[numbers=none,caption={micro-kernel descriptor 需要表达的信息}]
MatmulKernelDesc:
  isa: avx2 | avx512 | scalar
  input_dtype
  weight_format
  accumulator_dtype
  output_dtype
  tile_m, tile_n, tile_k
  alignment_bytes
  tail_policy
  quant_group_size
  requires_packed_weight
\end{lstlisting}

descriptor 的价值是把 kernel 能做什么说清楚。编译器 lowering 可以根据它选择 kernel；runtime profiler 可以把耗时归到具体 kernel；oracle 可以知道该用哪个误差阈值；测试可以枚举 tail、对齐和 dtype 组合。

\subsection{线程调度：小 kernel 不一定适合抢满所有核心}

CPU 推理的线程调度不能只写一句“并行 for”。prefill 和 decode 的形态不同，线程划分也不同。prefill 中 token 维度较大，GEMM/attention 更容易让多个核心保持忙碌；decode 中一次通常只有一个新 token，小 kernel 很多，过度并行会让 barrier、任务投递和 cache 竞争吞掉收益。

调度器至少要考虑：

\begin{itemize}
  \item 按输出行、输出列、token、head、layer 还是 batch 切分。
  \item 每个 task 的最小工作量，太小的 task 不值得进入线程池。
  \item workspace 和 output buffer 是否会 false sharing。
  \item 多层连续 segment 是否能减少 barrier 次数。
  \item 线程亲和和 NUMA 放置是否与权重内存一致。
  \item profiler 是否能分别记录 compute time、wait time 和 barrier time。
\end{itemize}

一个实用规则是：先让单线程 kernel 正确且可解释，再做粗粒度并行，最后才优化小任务调度。否则多线程版本慢了，很难判断是 kernel 慢、切分错、同步重，还是内存带宽已经满了。

\subsection{GPU 后端：共享语义合同，不共享 CPU 的物理假设}

GPU 后端和 CPU 后端应该共享 graph/tensor/quant/KV 的语义合同，但不共享所有物理布局。CPU 关心 cache line、SIMD 寄存器和线程池；GPU 关心显存连续性、coalesced load、shared memory、kernel launch、stream、同步和 host-device staging。

GPU 后端至少要有这些 C++20 RAII 边界：

\begin{itemize}
  \item \code{DeviceBuffer}：管理显存分配和释放，禁止裸指针泄漏生命周期。
  \item \code{Stream}：管理异步执行队列和事件同步。
  \item \code{Workspace}：为某些 kernel 或 vendor library 调用提供临时区。
  \item \code{KernelAdapter}：把 executable plan 的 descriptor 绑定到具体 kernel launch。
  \item \code{StagingBuffer}：处理 logits、token id 或小 metadata 的 host-device 传输。
\end{itemize}

GPU 优化要特别小心 decode 阶段。大 prefill 可能容易吃满 GPU，但 decode 单 token 可能被 launch overhead、同步、KV cache 读流或 logits 回传限制。若 profiler 只报告 GPU kernel 时间，不报告端到端 step latency，就会误判。

\subsection{后端优化的验收报告}

每次后端优化都应该产出一份小报告。它不一定长，但必须包含能复现和能判断的信息：

\begin{lstlisting}[numbers=none,caption={后端优化报告字段}]
optimization:
  name
  changed_contract
  affected_stage: prefill | decode | both
  backend: cpu | gpu

correctness:
  reference
  oracle_metric
  tolerance
  tested_shapes
  tested_dtypes

performance:
  baseline_commit_or_build
  hardware
  threads_or_device
  prompt_lengths
  context_lengths
  latency
  tokens_per_second
  memory_peak
  profiler_breakdown
\end{lstlisting}

这份报告就是工程质量的底线。没有 tested shapes，尾块可能没测；没有 tested dtypes，量化路径可能没测；没有 context length，KV cache 长上下文可能没测；没有 profiler breakdown，tokens/s 变快也不知道是哪部分变快。

\subsection{一次真实优化应该怎样设计实验}

以“把 decode 线性层从 fp16 权重改成 int4 weight-only”为例，不要直接改完跑一个 prompt。正确实验要先拆变量：

\begin{enumerate}
  \item 固定模型切片、prompt set、sampler seed、线程数或 GPU。
  \item 先验证 converter：原始权重坐标反量化与 reference 一致。
  \item 再验证单算子：同一输入下 fp16 reference 与 int4 kernel 输出误差。
  \item 再验证层输出：一个 decoder block 的 hidden drift。
  \item 再验证 logits：top-k 是否变化，最大误差和分位数是多少。
  \item 最后 benchmark：packing time、decode step latency、tokens/s、内存峰值。
\end{enumerate}

\begin{lstlisting}[numbers=none,caption={单次优化实验记录}]
Experiment:
  change: q4_weight_only_decode_linear
  fixed:
    model_fixture
    prompt_set
    sampler_seed
    backend
    thread_count
  correctness:
    byte_oracle_passed
    operator_error
    layer_error
    logits_topk_change
  performance:
    packing_time_ms
    decode_step_p50_p95
    tokens_per_second
    memory_peak_bytes
    bandwidth_estimate
\end{lstlisting}

这样实验失败时能定位。byte oracle 失败，查 nibble、scale、tail；operator 失败，查 kernel 地址流和累加；layer 失败，查 fusion、arena 或 residual；logits 失败，查误差是否被 softmax 放大；performance 没提升，查瓶颈是否根本不在权重带宽。

\subsection{A/B 路径和回滚条件}

后端优化必须可回滚。项目里应同时保留 reference、baseline 和 optimized plan，并让 runtime 能按配置选择。不是为了让用户长期使用慢路径，而是为了调试和线上回滚。

\begin{lstlisting}[numbers=none,caption={后端 plan 选择策略}]
BackendPlanSet:
  reference_plan
  baseline_plan
  optimized_plan
  debug_oracle_plan

Selection:
  if oracle_enabled:
    use debug_oracle_plan
  else if optimized_supported and gate_passed:
    use optimized_plan
  else:
    use baseline_plan
\end{lstlisting}

回滚条件要提前定义。例如：

\begin{itemize}
  \item logits top-k 变化率超过阈值；
  \item 固定 sampler 的 token mismatch 超过阈值；
  \item decode p95 latency 回归超过阈值；
  \item memory peak 超过容量预算；
  \item fallback 次数超过阈值；
  \item 某类 shape 或 context bucket 失败。
\end{itemize}

这不是保守，而是让优化可维护。没有 A/B 和回滚条件，项目会变成“某次优化之后好像快了，但后来没人敢动”。

\subsection{CPU perf counters 应该服务假设}

第二册讲过 cache、TLB、分支和指令吞吐。第三册用它们时要围绕假设，而不是堆计数器。比如你认为 int4 优化更快是因为少读权重，那么应观察：

\begin{itemize}
  \item LLC miss 或内存带宽是否下降；
  \item 指令数是否因为解包显著上升；
  \item cycles per token 是否下降；
  \item SIMD 指令比例是否符合预期；
  \item TLB miss 是否因 packed layout 改善或恶化。
\end{itemize}

\begin{lstlisting}[numbers=none,caption={CPU 后端性能假设记录}]
Hypothesis:
  q4 reduces weight bandwidth in decode GEMV

Counters:
  cycles
  instructions
  cache_misses
  dTLB_misses
  memory_bandwidth_estimate
  vector_instruction_count

Decision:
  if bytes_down but instructions_up too much:
    inspect unpack and scale path
  if cache_misses unchanged:
    bottleneck may not be weight stream
\end{lstlisting}

GPU 也同理。认为融合减少 launch，就要报告 launch count 和 total step latency；认为 KV layout 改善 coalescing，就要报告 memory throughput、load efficiency 或同类指标；认为 top-k on device 减少同步，就要报告 staging bytes 和 sampler wait。

\subsection{优化完成的定义}

一个后端优化完成，不是代码合并，也不是单个 benchmark 变快。最小完成定义是：

\begin{rubricbox}
\begin{itemize}
  \item descriptor 说明了新路径支持的 dtype、layout、shape、tail 和 quant profile。
  \item verifier 能拒绝不支持输入。
  \item reference/oracle 覆盖字节、单算子、层或 logits 中的相关层级。
  \item profiler 能拆分本次优化影响的时间和字节。
  \item benchmark 覆盖至少一个整齐形状和一个边界形状。
  \item 文档记录适用场景、收益、风险和回滚条件。
\end{itemize}
\end{rubricbox}

这套定义直接回答“学完能做成什么”。照这个标准推进，读者能做出一个有限但严肃的推理后端：支持明确模型范围、明确量化 profile、明确 CPU/GPU 能力、明确正确性阈值和性能报告。它未必马上追平成熟框架，但每个差距都能被定位。

\subsection{本节验收线}

读完本节，应该建立后端优化的基本纪律：

\begin{itemize}
  \item C++20 后端要分清 descriptor、RAII buffer、packed weight、kernel、scheduler、profiler 和 oracle。
  \item CPU 优化按 reference、scalar baseline、layout、packing、SIMD、线程、NUMA/TLB 的阶梯推进。
  \item SIMD micro-kernel 必须声明 dtype、layout、tile、tail、alignment 和数值误差合同。
  \item 线程调度要根据 prefill/decode 形态选择粒度，不能盲目填满线程池。
  \item GPU 后端共享语义合同，但用自己的 device buffer、stream、workspace 和 staging 边界。
  \item 每个优化都必须同时有 correctness report 和 performance report。
\end{itemize}

后续扩写具体 kernel 时，本节就是验收模板：没有合同、没有 oracle、没有 profiler、没有可复现 benchmark 的优化，都不能算完成。
